\documentclass[a4paper,12pt]{report}

% ---------- Packages ----------
\usepackage[utf8]{inputenc}
\usepackage[T1]{fontenc}
\usepackage[french]{babel}
\usepackage{graphicx}
\usepackage{hyperref}
\usepackage[top=2.5cm, bottom=2.5cm, left=3cm, right=2.5cm]{geometry}
\usepackage{fancyhdr}
\usepackage{titlesec}
\usepackage{setspace}
\usepackage{xcolor}
\usepackage{caption}
\usepackage{enumitem}
\usepackage{tocloft}

% ---------- Hyperlinks ----------
\hypersetup{
    colorlinks=true,
    linkcolor=teal,
    urlcolor=teal,
    citecolor=teal,
}

% ---------- Colors ----------
\definecolor{tealprimary}{RGB}{13,148,136}
\definecolor{tealdark}{RGB}{19,78,74}

% ---------- Page Style ----------
\onehalfspacing

\pagestyle{fancy}
\fancyhf{}
\fancyhead[L]{\textit{CampusEvents AI}}
\fancyhead[R]{\thepage}
\fancyfoot[C]{}
\renewcommand{\headrulewidth}{0.5pt}

\fancypagestyle{plain}{
    \fancyhf{}
    \fancyhead[L]{\textit{CampusEvents AI}}
    \fancyhead[R]{\thepage}
    \renewcommand{\headrulewidth}{0.5pt}
}

% ---------- Section Formatting ----------
\titleformat{\chapter}[display]
    {\Large\bfseries\color{tealdark}}
    {\chaptername\ \thechapter}
    {20pt}
    {\Huge}

\titleformat{\section}
    {\Large\bfseries\color{tealprimary}}
    {\thesection}
    {1em}
    {}

\titleformat{\subsection}
    {\large\bfseries\color{tealdark}}
    {\thesubsection}
    {1em}
    {}

% ---------- Table of Contents ----------
\renewcommand{\cfttoctitlefont}{\Huge\bfseries\color{tealdark}}
\renewcommand{\cftchapfont}{\bfseries\color{tealprimary}}
\renewcommand{\cftsecfont}{\color{black}}

% ---------- Document ----------
\begin{document}

% ============================================================
% COVER PAGE
% ============================================================
\begin{titlepage}
    \centering
    \vspace*{2cm}

    {\Huge\bfseries\color{tealdark} Rapport de Projet}\\[1.5cm]

    {\LARGE\bfseries\color{tealprimary} CampusEvents AI}\\[0.5cm]
    {\Large Application Mobile de Gestion d'\'{E}v\'{e}nements\\ avec Assistant IA}\\[2.5cm]

    \begin{minipage}{0.8\textwidth}
        \large
        \begin{minipage}[t]{0.45\textwidth}
            \raggedright
            \textbf{R\'{e}alis\'{e} par :}\\[0.3cm]
            {\Large OUALLALI Mohamed Amine}
        \end{minipage}
        \hfill
        \begin{minipage}[t]{0.45\textwidth}
            \raggedleft
            \textbf{Encadr\'{e} par :}\\[0.3cm]
            {\Large Hassan ZILI}
        \end{minipage}

        \vspace{1.5cm}
        \centering
        \href{https://github.com/mohamed-amine-dev/CampusEvents-AI.git}{\Large\color{tealprimary}\underline{https://github.com/mohamed-amine-dev/CampusEvents-AI.git}}

        \vspace{1.5cm}
        \centering
        \large
        \today
    \end{minipage}

    \vfill
\end{titlepage}

% ============================================================
% TABLE OF CONTENTS
% ============================================================
\tableofcontents
\newpage

% ============================================================
% INTRODUCTION
% ============================================================
\chapter*{Introduction}
\addcontentsline{toc}{chapter}{Introduction}

CampusEvents AI est une application mobile de gestion d'\'{e}v\'{e}nements campus, d\'{e}velopp\'{e}e dans le cadre d'un projet acad\'{e}mique. L'application vise \`{a} centraliser les \'{e}v\'{e}nements universitaires et \`{a} les rendre accessibles aux \'{e}tudiants et aux administrateurs, tout en int\'{e}grant un assistant intelligent bas\'{e} sur l'intelligence artificielle.

L'objectif principal est de fournir une plateforme intuitive o\`{u} les \'{e}tudiants peuvent d\'{e}couvrir des \'{e}v\'{e}nements, s'y inscrire, les sauvegarder en favoris, et interagir avec un assistant IA capable de r\'{e}pondre \`{a} leurs questions, de recommander des \'{e}v\'{e}nements pertinents et de les aider \`{a} planifier leur semaine. Les administrateurs disposent d'une interface compl\`{e}te pour cr\'{e}er, modifier et supprimer des \'{e}v\'{e}nements.

Ce rapport pr\'{e}sente l'architecture technique, les choix de conception, les fonctionnalit\'{e}s impl\'{e}ment\'{e}es et les perspectives d'\'{e}volution de l'application CampusEvents AI.

\newpage

% ============================================================
% ARCHITECTURE TECHNIQUE 
% ============================================================
\chapter*{Architecture Technique et Donn\'{e}es}
\addcontentsline{toc}{chapter}{Architecture Technique et Donn\'{e}es}

\section*{Stack Technologique}

CampusEvents AI est construite avec les technologies suivantes :

\begin{itemize}[leftmargin=2em]
    \item \textbf{React Native} (Expo SDK 56) : Framework cross-platform pour le d\'{e}veloppement d'applications mobiles natives
    \item \textbf{TypeScript} : Langage typ\'{e} pour une meilleure maintenabilit\'{e} du code
    \item \textbf{Expo Router} : Syst\`{e}me de navigation bas\'{e} sur le syst\`{e}me de fichiers inspir\'{e} de Next.js
    \item \textbf{SQLite} : Base de donn\'{e}es locale embarqu\'{e}e
    \item \textbf{Groq API} : Service d'inf\'{e}rence LLM pour l'assistant IA (mod\`{e}le Llama 3.3 70B)
    \item \textbf{AsyncStorage} : Stockage persistant pour les sessions utilisateur et le cache
\end{itemize}

\section*{Structure du Projet}

L'application suit une architecture modulaire organis\'{e}e autour du routing bas\'{e} fichiers :

\begin{itemize}[leftmargin=2em]
    \item \texttt{app/} : Pages et navigation via Expo Router (groupes \'{e}tudiants et administrateurs)
    \item \texttt{components/ui/} : 9 composants r\'{e}utilisables (Button, Card, TextField, SearchBar, Chip, Badge, Header, Screen, EmptyState)
    \item \texttt{context/} : Contextes React (authentification et th\`{e}me)
    \item \texttt{database/} : Couche d'acc\`{e}s aux donn\'{e}es avec moteur SQL personnalis\'{e}
    \item \texttt{services/} : Int\'{e}gration de l'API Groq
    \item \texttt{constants/} : Tokens de design, couleurs, donn\'{e}es de d\'{e}monstration
    \item \texttt{types/} : D\'{e}finitions TypeScript
\end{itemize}

\section*{Couche de Donn\'{e}es}

La base de donn\'{e}es comprend quatre tables principales :

\begin{itemize}[leftmargin=2em]
    \item \textbf{events} : stocke les \'{e}v\'{e}nements (titre, description, cat\'{e}gorie, date, lieu, organisateur, capacit\'{e}, image)
    \item \textbf{registrations} : enregistre les inscriptions des utilisateurs aux \'{e}v\'{e}nements
    \item \textbf{favorites} : sauvegarde les favoris des utilisateurs
    \item \textbf{llm\_results} : met en cache les r\'{e}ponses de l'assistant IA
\end{itemize}

La couche \texttt{database/db.web.ts} impl\'{e}mente un moteur SQL complet en m\'{e}moire avec persistance AsyncStorage. Ce moteur analyse les requ\^{e}tes SQL via des expressions r\'{e}guli\`{e}res et supporte les op\'{e}rations CRUD compl\`{e}tes avec support des clauses WHERE, ORDER BY, LIMIT, INSERT OR IGNORE et des expressions de mise \`{a} jour.

Pour la plateforme native, \texttt{database/db.native.ts} utilise l'API synchrone d'\texttt{expo-sqlite}. Le s\'{e}lecteur \texttt{database/db.ts} choisit automatiquement l'impl\'{e}mentation appropri\'{e}e selon la plateforme d\'{e}tect\'{e}e.

\section*{Donn\'{e}es de D\'{e}monstration}

L'application est livr\'{e}e avec 12 \'{e}v\'{e}nements pr\'{e}configur\'{e}s couvrant cinq cat\'{e}gories : Workshop, Talk, Club, Exam et Other. Ces \'{e}v\'{e}nements sont g\'{e}n\'{e}r\'{e}s avec des dates dynamiques relatives \`{a} la date courante et incluent des images Unsplash th\'{e}matiques. Deux comptes de d\'{e}monstration sont pr\'{e}configur\'{e}s :

\begin{itemize}[leftmargin=2em]
    \item \texttt{admin@campus.ma} / \texttt{admin123} (r\^{o}le administrateur)
    \item \texttt{etudiant@campus.ma} / \texttt{etudiant123} (r\^{o}le \'{e}tudiant)
\end{itemize}

\newpage

% ============================================================
% FONCTIONNALITÉS
% ============================================================
\chapter*{Fonctionnalit\'{e}s}
\addcontentsline{toc}{chapter}{Fonctionnalit\'{e}s}

\section*{Authentification et Navigation}

Le syst\`{e}me d'authentification repose sur un \texttt{AuthContext} React qui persiste la session utilisateur dans AsyncStorage. La page de connexion propose un s\'{e}lecteur de r\^{o}le (Admin/\'{E}tudiant) qui pr\'{e}-remplit automatiquement les identifiants correspondants.

La navigation est enti\`{e}rement g\'{e}r\'{e}e par Expo Router avec deux groupes d'\'{e}crans distincts :

\begin{itemize}[leftmargin=2em]
    \item \textbf{Groupe \'{e}tudiant} \texttt{(student)} : navigation par onglets avec 5 \'{e}crans (Catalogue, Favoris, Assistant IA, Inscriptions, Profil)
    \item \textbf{Groupe administrateur} \texttt{(admin)} : navigation par pile avec 3 \'{e}crans (Liste des \'{e}v\'{e}nements, Cr\'{e}ation, Modification)
\end{itemize}

Une garde d'authentification dans le layout racine redirige automatiquement vers la page de connexion si l'utilisateur n'est pas connect\'{e}.

\section*{Fonctionnalit\'{e}s \'{E}tudiant}

\subsection*{Catalogue}
Le catalogue affiche les \'{e}v\'{e}nements dans une grille \`{a} deux colonnes avec recherche textuelle et filtres par cat\'{e}gorie. Chaque carte d'\'{e}v\'{e}nement pr\'{e}sente une image hero, un badge de date, la jauge de capacit\'{e}, le nom et la localisation. Un bouton c{\oe}ur permet d'ajouter ou retirer un favori instantan\'{e}ment.

\subsection*{Favoris et Inscriptions}
Les \'{e}crans Favoris et Inscriptions listent respectivement les \'{e}v\'{e}nements favoris et les inscriptions confirm\'{e}es de l'utilisateur. L'inscription se fait depuis la page de d\'{e}tail d'un \'{e}v\'{e}nement, avec une confirmation visuelle et une mise \`{a} jour en temps r\'{e}el du compteur de participants.

\subsection*{D\'{e}tail d'\'{E}v\'{e}nement}
La page de d\'{e}tail pr\'{e}sente une section hero avec l'image, les informations cl\'{e}s (lieu, organisateur, capacit\'{e} avec barre de progression), la description compl\`{e}te et les tags. Un bouton principal permet de s'inscrire ou d'annuler l'inscription.

\subsection*{Profil et Th\`{e}me}
L'utilisateur peut modifier son nom d'affichage et basculer entre les th\`{e}mes clair et sombre. Le th\`{e}me est persist\'{e} dans AsyncStorage et appliqu\'{e} globalement via un \texttt{ThemeContext}.

\section*{Fonctionnalit\'{e}s Administrateur}

Les administrateurs disposent d'un panneau complet de gestion des \'{e}v\'{e}nements :

\begin{itemize}[leftmargin=2em]
    \item \textbf{Liste} : visualisation de tous les \'{e}v\'{e}nements avec indicateurs de capacit\'{e} et suppression
    \item \textbf{Cr\'{e}ation} : formulaire complet avec validation des champs requis, s\'{e}lecteur de cat\'{e}gorie et date
    \item \textbf{Modification} : pr\'{e}-remplissage du formulaire avec les donn\'{e}es existantes
    \item \textbf{Suppression} : confirmation avant suppression avec cascade dans la base de donn\'{e}es
\end{itemize}

\newpage

% ============================================================
% ASSISTANT IA
% ============================================================
\chapter*{Assistant IA}
\addcontentsline{toc}{chapter}{Assistant IA}

\section*{Int\'{e}gration Groq}

L'assistant IA utilise l'API Groq avec le mod\`{e}le \texttt{Llama 3.3 70B Versatile}. La cl\'{e} API est configur\'{e}e via une variable d'environnement (\texttt{EXPO\_PUBLIC\_GROQ\_API\_KEY}) et accessible au client gr\^{a}ce au pr\'{e}fixe EXPO\_PUBLIC d'Expo.

\section*{Modes de l'Assistant}

L'assistant propose quatre modes sp\'{e}cialis\'{e}s, accessibles via des onglets :

\begin{itemize}[leftmargin=2em]
    \item \textbf{Recherche} : interpr\`{e}te les requ\^{e}tes en langage naturel et retourne les identifiants des \'{e}v\'{e}nements correspondants sous forme de tableau JSON
    \item \textbf{Recommandations} : sugg\`{e}re des \'{e}v\'{e}nements personnalis\'{e}s en fonction des favoris et inscriptions de l'utilisateur
    \item \textbf{Planning} : aide l'utilisateur \`{a} organiser sa semaine en analysant les \'{e}v\'{e}nements \`{a} venir dans les 7 prochains jours
    \item \textbf{Questions} : r\'{e}pond aux questions g\'{e}n\'{e}rales sur les \'{e}v\'{e}nements du campus en fran\c{c}ais
\end{itemize}

\section*{Cache et Performance}

Pour optimiser les performances et r\'{e}duire les appels API, les r\'{e}sultats de l'assistant sont mis en cache dans la table \texttt{llm\_results}. Une requ\^{e}te identique (m\^{e}me utilisateur, m\^{e}me type, m\^{e}me texte d'entr\'{e}e) renvoie le r\'{e}sultat en cache sans solliciter l'API. Cette approche permet \'{e}galement un historique des interactions.

\section*{Ing\'{e}nierie des Prompts}

Chaque mode utilise un prompt syst\`{e}me sp\'{e}cialis\'{e} avec un contexte riche :

\begin{itemize}[leftmargin=2em]
    \item \textbf{Search} : renvoie exclusivement un tableau JSON d'IDs d'\'{e}v\'{e}nements correspondant \`{a} la requ\^{e}te
    \item \textbf{Recommendation} : contexte enrichi avec les favoris et inscriptions de l'utilisateur pour des suggestions pertinentes
    \item \textbf{Planning} : filtre temporel sur la semaine \`{a} venir avec d\'{e}tails horaires et lieux
    \item \textbf{QA} : base de connaissances compl\`{e}te de tous les \'{e}v\'{e}nements pour des r\'{e}ponses pr\'{e}cises
\end{itemize}

La temp\'{e}rature est fix\'{e}e \`{a} 0.3 pour des r\'{e}ponses pr\'{e}cises et reproductibles, et le nombre maximal de tokens est limit\'{e} \`{a} 2000.

\newpage

% ============================================================
% CONCLUSION
% ============================================================
\chapter*{Conclusion}
\addcontentsline{toc}{chapter}{Conclusion}

CampusEvents AI est une application compl\`{e}te qui d\'{e}montre l'int\'{e}gration r\'{e}ussie de technologies modernes de d\'{e}veloppement mobile avec l'intelligence artificielle. L'application offre une exp\'{e}rience utilisateur fluide avec une interface soign\'{e}e, un th\`{e}me adaptatif, et des fonctionnalit\'{e}s avanc\'{e}es comme l'assistant IA.

Les principaux acquis de ce projet incluent :

\begin{itemize}[leftmargin=2em]
    \item La ma\^{i}trise du framework React Native et de l'\'{e}cosyst\`{e}me Expo
    \item L'impl\'{e}mentation d'une couche de donn\'{e}es multi-plateforme (web et natif)
    \item L'int\'{e}gration d'un LLM via une API REST avec gestion de cache
    \item La conception d'une architecture modulaire et maintenable
    \item La mise en place d'un syst\`{e}me de design coh\'{e}rent avec support du mode sombre
\end{itemize}

\vspace{0.5cm}
Parmi les perspectives d'am\'{e}lioration :
\begin{itemize}[leftmargin=2em]
    \item Synchronisation avec un backend distant pour le partage multi-appareils
    \item Notifications push pour les rappels d'\'{e}v\'{e}nements
    \item Syst\`{e}me de feedback et notation des \'{e}v\'{e}nements
    \item Support multilingue (anglais, arabe)
    \item Am\'{e}lioration des prompts de l'assistant avec contexte conversationnel
\end{itemize}

\vspace{1cm}
\centering
\large
\href{https://github.com/mohamed-amine-dev/CampusEvents-AI.git}{\color{tealprimary}\underline{https://github.com/mohamed-amine-dev/CampusEvents-AI.git}}

\newpage

% ============================================================
% ANNEXE : CAPTURES D'ÉCRAN
% ============================================================
\chapter*{Annexe : Captures d'\'{E}cran}
\addcontentsline{toc}{chapter}{Annexe : Captures d'\'{E}cran}

\begin{figure}[h]
    \centering
    \includegraphics[width=0.7\textwidth]{img1.jpeg}
    \caption{Page de connexion avec s\'{e}lecteur de r\^{o}le}
\end{figure}
\newpage

\begin{figure}[h]
    \centering
    \includegraphics[width=0.7\textwidth]{img2.jpeg}
    \caption{Catalogue d'\'{e}v\'{e}nements en grille 2 colonnes}
\end{figure}
\newpage

\begin{figure}[h]
    \centering
    \includegraphics[width=0.7\textwidth]{img3.jpeg}
    \caption{Filtres par cat\'{e}gorie et recherche textuelle}
\end{figure}
\newpage

\begin{figure}[h]
    \centering
    
\includegraphics[width=0.7\textwidth]{img4.jpeg}
    \caption{D\'{e}tail d'un \'{e}v\'{e}nement avec hero image}
\end{figure}
\newpage

\begin{figure}[h]
    \centering
    \includegraphics[width=0.7\textwidth]{img5.jpeg}
    \caption{Inscription et barre de capacit\'{e}}
\end{figure}
\newpage

\begin{figure}[h]
    \centering
    
\includegraphics[width=0.7\textwidth]{img6.jpeg}
    \caption{\'{E}cran des favoris}
\end{figure}
\newpage

\begin{figure}[h]
    \centering
    
\includegraphics[width=0.7\textwidth]{img7.jpeg}
    \caption{\'{E}cran des inscriptions avec option d'annulation}
\end{figure}
\newpage

\begin{figure}[h]
    \centering
    \includegraphics[width=0.7\textwidth]{img8.jpeg}
    \caption{Assistant IA en mode Recommandations}
\end{figure}
\newpage

\begin{figure}[h]
    \centering
    \includegraphics[width=0.7\textwidth]{img9.jpeg}
    \caption{Gestion des \'{e}v\'{e}nements (interface administrateur)}
\end{figure}
\newpage

\begin{figure}[h]
    \centering
    \includegraphics[width=0.7\textwidth]{img10.jpeg}
    \caption{Formulaire de cr\'{e}ation d'\'{e}v\'{e}nement}
\end{figure}
\newpage

\begin{figure}[h]
    \centering
    \includegraphics[width=0.7\textwidth]{img11.jpeg}
    \caption{Profil utilisateur et s\'{e}lecteur de th\`{e}me}
\end{figure}

\end{document}
